\documentclass[12pt]{article}
\usepackage{cmap} % Маст-хэв для корректной работы с кодировками и поиска в PDF
\usepackage[utf8]{inputenc}
\usepackage[T2A]{fontenc}
\usepackage[russian]{babel}
\usepackage{amsmath}
\usepackage{amssymb}
\usepackage{geometry}
\usepackage{enumitem}
\usepackage{hyperref}
\usepackage{xcolor}
\usepackage{titlesec}
\usepackage{parskip}
\usepackage{listings}
\usepackage[expansion=false]{microtype} % Отключаем expansion, чтобы убрать ошибку шрифтов

% Настройка полей
\geometry{a4paper, margin=2cm}

% Цвета
\definecolor{primary}{RGB}{0, 102, 204}
\definecolor{secondary}{RGB}{100, 100, 100}
\definecolor{summarybg}{RGB}{245, 245, 250}
\definecolor{extra_info}{RGB}{0, 102, 51} % Спокойный зеленый для дополнений

% Настройка listings
\lstset{
	language=Python,
	basicstyle=\small\ttfamily,
	keywordstyle=\color{blue},
	commentstyle=\color{gray},
	stringstyle=\color{orange},
	breaklines=true,
	frame=single,
	backgroundcolor=\color{summarybg}
}

% Стиль заголовков
\titleformat{\section}{\color{primary}\normalfont\Large\bfseries}{\thesection.}{0.5em}{}
\titleformat{\subsection}{\color{primary}\normalfont\large\bfseries}{}{0em}{}

% Настройка списков
\setlist[itemize]{noitemsep, topsep=0pt}
\setlist[enumerate]{font=\color{primary}\bfseries}

\begin{document}
	
	\begin{center}
		{\huge\bfseries\color{primary} Конспект для подготовки к собеседованию} \\
		\vspace{0.5em}
		{\large Data Science (на основе видео и ML Handbook Яндекса)}
	\end{center}
	
	\vspace{1em}
	\noindent \textbf{Вот расширенный и глубоко проработанный конспект для подготовки к собеседованию. Я значительно детализировал цитаты из видео, добавил экспертные «фишки», о которых говорит спикер, и обогатил теоретическую часть математическими выкладками и строгими концепциями из ML Handbook.} [cite: 1, 2]
	
	\hrule
	\vspace{1.5em}
	
	\begin{enumerate}[leftmargin=*]
		\item \textbf{Как работает логистическая регрессия.}
		\begin{description}[leftmargin=1em]
			\item[\textcolor{primary}{Цитата из видео:}] «У нас транспонированный вектор весов признаков умножается на наш объект. Мы пропускаем это всё через сигмоиду и получаем вероятность принадлежности классу. <...> Фишечка: расскажите почему именно логлос. То, что это приходит из теории вероятности, из максимизации правдоподобия. Расскажите то, что он выпукл и вы будете всегда сходиться к одному решению. <...> На больших ошибках предсказаний он даёт большой штраф». [cite: 3, 4, 5, 6]
			\item[\textcolor{primary}{Дополнение из учебника:}] Логистическая регрессия предсказывает логит (log odds) — логарифм отношения вероятностей:
			\[ \log\left(\frac{p}{1-p}\right) \]
			Чтобы перевести линейный ответ $\langle w, x_i \rangle$ в вероятность, используется сигмоида:
			\[ \sigma(z) = \frac{1}{1 + e^{-z}} \]
			Оптимизация LogLoss — это следствие метода максимального правдоподобия для распределения Бернулли:
			\[ p(y|X,w) = \prod p_i^{y_i}(1-p_i)^{1-y_i} \]
			
			Разделяющая поверхность классификатора задается уравнением $\langle w,x \rangle = 0$, что геометрически является гиперплоскостью. [cite: 7]
			\item[\textcolor{primary}{Резюме:}] Это линейный классификатор, обучаемый через максимизацию правдоподобия (минимизацию строго выпуклого LogLoss). Модель предсказывает логарифм отношения шансов, который через сигмоиду трансформируется в вероятность, а граница классов образует гиперплоскость. [cite: 8, 9]
		\end{description}
		
		\textcolor{extra_info}{\textbf{Дополнительный пример:} В задачах банковского скоринга лог-лос критически важен, так как он сильно штрафует модель за ложную уверенность. Если модель предсказывает вероятность 0.99 для «плохого» заемщика, штраф будет намного выше, чем при предсказании 0.6.}
		
		\item \textbf{Что такое переобучение и как его замечать.}
		\begin{description}[leftmargin=1em]
			\item[\textcolor{primary}{Цитата из видео:}] «Когда наша модель слишком сильно подстроилась под данные: у нас есть неплохое качество на трейне, но на тесте качество слишком плохое. <...> Более продвинутый ответ — это использовать кросс-валидацию, посмотреть, что модель подстроилась под различные фолды. <...> У каждой модели есть своя специфика борьбы: у деревьев это высота листа, у нейронных сетей — дропаут». [cite: 10, 11, 12]
			\item[\textcolor{primary}{Дополнение из учебника:}] В терминах теории ML переобучение описывается через Bias-Variance Decomposition. [cite: 13] Переобученная модель имеет высокий разброс (Variance) — она избыточно чувствительна к случайным флуктуациям и запоминает шум обучающей выборки вместо реального сигнала. [cite: 14]
			\item[\textcolor{primary}{Резюме:}] Переобучение (Overfitting) — это состояние модели с высоким Variance, когда она заучивает шум. [cite: 15] Детектируется через разрыв метрик между train и test/фолдами кросс-валидации. Лечится регуляризацией, зависящей от типа алгоритма. [cite: 16]
		\end{description}
		
		\item \textbf{Фишки: data augmentation, SMOTE, шум.}
		\begin{description}[leftmargin=1em]
			\item[\textcolor{primary}{Цитата из видео:}] «Круто работать с переобучением через данные: можно добавить шум, можно сделать больше новых данных... использовать SMOTE, чтобы разнообразить ваш датасет. Это всё очень хорошие методы борьбы с переобучением». [cite: 17]
			\item[\textcolor{primary}{Дополнение из учебника:}] Аугментация (внесение изменений в данные) — важнейший инструмент регуляризации в глубинном обучении. [cite: 18] Например, поворот картинки самолета на 10 градусов создает новый для сети объект, но с той же меткой. [cite: 19] Это позволяет «дать модели понять, какие преобразования над данными являются допустимыми». [cite: 20] Главное ограничение: аугментированные данные должны принадлежать к той же генеральной совокупности. [cite: 21]
			\item[\textcolor{primary}{Резюме:}] Аугментация и искусственное зашумление данных — это мощная форма регуляризации, расширяющая выборку и обучающая модель инвариантности к допустимым в реальном мире искажениям (повороты, шумы). [cite: 22]
		\end{description}
		
		\item \textbf{Почему L1 регуляризация зануляет признаки.}
		\begin{description}[leftmargin=1em]
			\item[\textcolor{primary}{Цитата из видео:}] «Интуитивное понимание: гиперплоскость L1 регуляризации — это октаэдр, а L2 — сфера. И октаэдр сводится к нулям. <...> Более крутое доказательство — это через субдифференциалы: мы доопределяем производную по модулю. У L2 идут большие штрафы при больших весах, а для модуля градиент более равномерный. Поэтому веса равномерно отнимаются и сходятся к нулю». [cite: 23, 24, 25]
			\item[\textcolor{primary}{Дополнение из учебника:}] При L1-регуляризации (Lasso) штраф имеет вид:
			\[ \lambda \sum |w_j| \] [cite: 26]
			Линии уровня L1-нормы — это ромбы (в 2D) или многогранники с острыми вершинами, лежащими строго на осях координат. [cite: 27] Линии уровня основной функции потерь в точке оптимума касаются именно этих углов на осях, обнуляя часть координат. [cite: 28] Оптимизация L1 сложна из-за недифференцируемости модуля в нуле, для чего применяются проксимальные методы градиентного спуска. [cite: 29]
			\item[\textcolor{primary}{Резюме:}] Геометрия L1-штрафа (ромбы с острыми углами на осях) и свойства субградиента приводят к тому, что точка касания с функцией потерь попадает на оси координат, создавая разреженные (зануленные) вектора весов, что работает как встроенный feature selection. [cite: 30]
		\end{description}
		
		\item \textbf{Как устроено дерево решений.}
		\begin{description}[leftmargin=1em]
			\item[\textcolor{primary}{Цитата из видео:}] «Это бинарное дерево. Для категорий есть ответы да/нет. Для вещественных признаков мы разбиваем их на бины, проходим по порядку и смотрим: если признак меньше значения — идет влево, больше — вправо. Мы хотим уменьшать кроссэнтропию или увеличивать Information Gain (чтобы в листьях была как можно меньше дисперсия)». [cite: 31, 32, 33]
			\item[\textcolor{primary}{Дополнение из учебника:}] Дерево задает кусочно-постоянную аппроксимацию. Внутренняя вершина содержит предикат:
			\[ B_v(x, j, t) = [x_j \le t] \] [cite: 34]
			Жадный алгоритм на каждом шаге максимизирует критерий ветвления:
			\[ Branch = |X_m|H(X_m) - |X_l|H(X_l) - |X_r|H(X_r) \]
			Информативность (impurity) $H(X)$ зависит от задачи: [cite: 35]
			\begin{itemize}
				\item для регрессии (MSE) это дисперсия ответов: $\frac{1}{|X|} \sum(y_i - \bar{y})^2$
				\item для классификации — критерий Джини: $\sum p_k(1-p_k)$
				\item или энтропия Шеннона: $-\sum p_k \log p_k$
			\end{itemize}
			\item[\textcolor{primary}{Резюме:}] Дерево строится жадно: пространство признаков итеративно делится гиперплоскостями, параллельными осям. [cite: 36] Сплит выбирается так, чтобы максимизировать падение неопределенности (дисперсии таргета или энтропии классов) в дочерних узлах. [cite: 37]
		\end{description}
		
		\item \textbf{Переобучение в деревьях, прунинг.}
		\begin{description}[leftmargin=1em]
			\item[\textcolor{primary}{Цитата из видео:}] «Дерево можно ограничить не с помощью кросс-валидации, а при построении: добавить в лосс наказание плюс $\alpha$ умножить на высоту дерева. Расскажите про прунинг. Можно требовать минимальное количество элементов в листе при разбитии». [cite: 38, 39]
			\item[\textcolor{primary}{Дополнение из учебника:}] Деревья — чемпионы переобучения, они могут идеально приблизить обучающую выборку (до 1 объекта в листе). [cite: 40] Регуляризация (early stopping/пре-прунинг) останавливает ветвление при достижении лимитов глубины или размера листа. [cite: 41] Прунинг (стрижка) работает иначе: дерево строится жадно до конца, а затем вершины удаляются с проверкой качества на отложенной выборке. [cite: 42]
			\item[\textcolor{primary}{Резюме:}] Деревья легко запоминают шум. С этим борются либо останавливая построение (early stopping), либо отсекая лишние ветви уже построенного глубокого дерева на отложенной выборке (прунинг). [cite: 43]
		\end{description}
		
		\item \textbf{Как работает случайный лес (Random Forest).}
		\begin{description}[leftmargin=1em]
			\item[\textcolor{primary}{Цитата из видео:}] «Деревья строятся параллельно и независимо друг от друга. Используется бутстрап объектов. Важная фишка: признаки выбираются бутстрапом *на каждом разбиении* (в каждом узле), а не для всего дерева в начале. Лес нужен, чтобы уменьшать дисперсию ошибки (variance), поэтому нам нужны как можно более глубокие, переобученные деревья с минимальным смещением». [cite: 44, 45, 46]
			\item[\textcolor{primary}{Дополнение из учебника:}] В основе лежит метод случайных подпространств и бэггинг. В каждой вершине отбираются случайные $n < N$ признаков. [cite: 47] Если базовые глубокие деревья не скоррелированы, то дисперсия ансамбля из $k$ деревьев падает в $k$ раз:
			\[ \mathbb{V}[a(x,X)] = \frac{1}{k}\mathbb{V}_X b(x,X) \]
			Смещение (bias) леса при этом остается равным смещению одного дерева. [cite: 48]
			\item[\textcolor{primary}{Резюме:}] Случайный лес — параллельный ансамбль глубоких переобученных деревьев (с низким bias). [cite: 49] Рандомизация по объектам (бутстрап) и признакам делает деревья независимыми, что при усреднении радикально снижает общую дисперсию (variance). [cite: 50]
		\end{description}
		
		\item \textbf{Чем отличается бустинг от Random Forest.}
		\begin{description}[leftmargin=1em]
			\item[\textcolor{primary}{Цитата из видео:}] «Бустинг строится последовательно, каждая модель учится на ошибку предыдущей. Если выкинуть первую модель, бустинг сломается, а случайный лес — нет. Бустинг уменьшает смещение (bias), поэтому мы хотим, чтобы дисперсия каждой базовой модели была как можно меньше — используются неглубокие деревья (от 3 до 8 уровней)». [cite: 51, 52, 53]
			\item[\textcolor{primary}{Дополнение из учебника:}] Градиентный бустинг строит композицию аддитивно. Каждое следующее дерево настраивается на минимизацию антиградиента функции потерь композиции всех прошлых шагов. [cite: 54] В отличие от леса, бустинг требует простых базовых алгоритмов для постепенного снижения Bias, но подвержен риску переобучения. [cite: 55]
			\item[\textcolor{primary}{Резюме:}] Лес снижает Variance, усредняя параллельные глубокие деревья. Бустинг снижает Bias, строя последовательность мелких деревьев, каждое из которых корректирует остаточную ошибку предыдущих. [cite: 56]
		\end{description}
		
		\item \textbf{Метрики регрессии и классификации.}
		\begin{description}[leftmargin=1em]
			\item[\textcolor{primary}{Цитата из видео:}] «Регрессия: MAE, MAPE, MSE, RMSE, SMAPE. Важно говорить, что они значат (MSE подстраивается под выбросы, MAPE нужна для понимания порядка). Классификация: Accuracy, Precision, Recall, F1, F-beta, ROC-AUC, PR-кривая». [cite: 57, 58]
			\item[\textcolor{primary}{Дополнение из учебника:}] MSE математически приближает матожидание целевой переменной. [cite: 59] MAE приближает медиану и значительно меньше штрафует за грубые выбросы. [cite: 60] Accuracy неинформативна при дисбалансе классов. [cite: 61]
			\item[\textcolor{primary}{Резюме:}] Для регрессии метрики выбираются исходя из отношения к выбросам и необходимости относительных оценок. [cite: 62] В классификации при дисбалансе фокус смещается на метрики полноты, точности и ранжирования. [cite: 63]
		\end{description}
		
		\item \textbf{Что такое ROC-AUC и как его интерпретировать.}
		\begin{description}[leftmargin=1em]
			\item[\textcolor{primary}{Цитата из видео:}] «ROC-AUC показывает: если мы берём случайный объект из класса 1 и случайный объект из класса 0, наша модель их правильно отранжировала. Вопрос-ловушка: как изменится ROC-AUC, если предсказания 0.8 и 0.3 заменить на 0.9 и 0.7? Ответ: никак, нам не важна сама вероятность, важно лишь расположение». [cite: 64, 65, 66]
			\item[\textcolor{primary}{Дополнение из учебника:}] Это площадь под кривой (Area Under Curve), отражающей зависимость True Positive Rate от False Positive Rate при всевозможных порогах. [cite: 67]
			\item[\textcolor{primary}{Резюме:}] ROC-AUC — это вероятность правильного попарного ранжирования: случайно взятый положительный объект получит скор выше, чем случайно взятый отрицательный. [cite: 68]
		\end{description}
		
		\item \textbf{Precision — как объяснить на собеседовании.}
		\begin{description}[leftmargin=1em]
			\item[\textcolor{primary}{Цитата из видео:}] «Точность. Нарисуйте для себя Confusion Matrix. Пример: выбор места для постройки дома. Нам не страшно пропустить хорошие места, но очень важно, чтобы то место, которое мы определили как "хорошее", гарантированно не было подвержено катаклизмам». [cite: 69, 70, 71]
			\item[\textcolor{primary}{Дополнение из учебника:}] Доля истинно положительных среди всех объектов, названных классификатором положительными. [cite: 72] Служит для контроля ошибок первого рода.
			\item[\textcolor{primary}{Резюме:}] Точность измеряет "надежность" положительных срабатываний модели. [cite: 73] Метрика критична там, где цена ложной тревоги катастрофически высока. [cite: 74]
		\end{description}
		
		\item \textbf{Recall и его применение в медицине.}
		\begin{description}[leftmargin=1em]
			\item[\textcolor{primary}{Цитата из видео:}] «Пример: в медицине нам очень важно не пропустить больного. Для нас не так страшно, если мы здорового отправим на обследование, главное — человек выживет». [cite: 75, 76]
			\item[\textcolor{primary}{Дополнение из учебника:}] Полнота (Recall) характеризует способность модели находить целевые объекты из всего множества реально существующих положительных примеров. [cite: 77] Фокусируется на минимизации ошибок второго рода.
			\item[\textcolor{primary}{Резюме:}] Полнота измеряет "охват" класса. [cite: 78] Ее максимизируют, когда пропуск события обходится намного дороже, чем ложное срабатывание. [cite: 79]
		\end{description}
		
		\item \textbf{Связь Recall и ROC-AUC.}
		\begin{description}[leftmargin=1em]
			\item[\textcolor{primary}{Цитата из видео:}] «Фишечка: Recall на самом деле используется в ROC-AUC. Осями там выступают True Positive Rate (это Recall для класса 1) и False Positive Rate». [cite: 80, 81]
			\item[\textcolor{primary}{Дополнение из учебника:}] Кривая строится в координатах (FPR, TPR) при изменении порога уверенности. TPR математически тождественна метрике Recall. [cite: 82]
			\item[\textcolor{primary}{Резюме:}] ROC-кривая буквально состоит из метрик полноты: она показывает, как меняется полнота нахождения целевого класса ценой потери полноты правильного нахождения отрицательного класса. [cite: 83]
		\end{description}
		
		\item \textbf{Зачем нужна PR-кривая и в чём её отличие от ROC-AUC.}
		\begin{description}[leftmargin=1em]
			\item[\textcolor{primary}{Цитата из видео:}] «При суперсильном дисбалансе классов PR-кривая показывает себя лучше. ROC-AUC более оптимистичен, а PR-кривая делает акцент именно на положительных таргетах». [cite: 84, 85]
			\item[\textcolor{primary}{Дополнение из учебника:}] PR-кривая концентрируется исключительно на качестве предсказания меньшего (положительного) класса, игнорируя успехи в нахождении True Negatives. [cite: 86]
			\item[\textcolor{primary}{Резюме:}] В условиях жесткого дисбаланса ROC-AUC дает обманчиво высокие оценки. [cite: 87] PR-кривая фокусируется только на миноритарном классе, давая объективную картину.
		\end{description}
		
		\item \textbf{Подбор гиперпараметров (GridSearch, RandomSearch, Optuna).}
		\begin{description}[leftmargin=1em]
			\item[\textcolor{primary}{Цитата из видео:}] «GridSearch не подходит, если большой набор значений. RandomSearch берет случайные значения. Байесовская оптимизация строит распределение потенциальной награды». [cite: 88, 89]
			\item[\textcolor{primary}{Дополнение из учебника:}] GridSearch страдает от размерности. RandomSearch эффективнее находит оптимум (60 итераций дают вероятность 95\% попасть в топ-5\% лучших решений). [cite: 90] Байесовская оптимизация учится на истории запусков. [cite: 91]
			\item[\textcolor{primary}{Резюме:}] Перебор по сетке слишком долгий. Случайный поиск работает лучше за счет исследования пространства. [cite: 92] Индустриальный стандарт — Байесовская оптимизация (Optuna). [cite: 93]
		\end{description}
		
		\item \textbf{Как оценить важность признаков (feature importance) в модели.}
		\begin{description}[leftmargin=1em]
			\item[\textcolor{primary}{Цитата из видео:}] «В бизнесе важно показывать, почему ML работает. Линейные: коэффициенты. Деревья: индекс Джини. Permutation Importance: перемешали значения фичи, посмотрели на качество. SHAP values — на основе теории игр». [cite: 94, 95, 96]
			\item[\textcolor{primary}{Дополнение из учебника:}] В линейной модели веса имеют смысл только при отсутствии мультиколлинеарности. [cite: 97, 98]
			\item[\textcolor{primary}{Резюме:}] Базовые методы — веса или InfoGain. Универсальные методы — Permutation и SHAP. [cite: 99] SHAP математически самый точный, но вычислительно тяжелый.
		\end{description}
		
		\item \textbf{Как работать с категориальными фичами (и чем хорош CatBoost).}
		\begin{description}[leftmargin=1em]
			\item[\textcolor{primary}{Цитата из видео:}] «Базовые: Label Encoding, One-Hot, Frequency, Target Encoding. CatBoost хорош тем, что сам берет пары категорий и использует Ordered boosting». [cite: 100, 101, 102]
			\item[\textcolor{primary}{Дополнение из учебника:}] При классическом One-Hot возникает риск мультиколлинеарности (dummy variable trap). [cite: 103]
			\item[\textcolor{primary}{Резюме:}] Простые методы кодируют категории числами или векторами. Продвинутые модели (CatBoost) генерируют комбинации и используют динамическое таргет-кодирование. [cite: 104]
		\end{description}
		
		\item \textbf{В чём разница между \texttt{==} и \texttt{is} в Python.}
		\begin{description}[leftmargin=1em]
			\item[\textcolor{primary}{Цитата из видео:}] «Базовый ответ: \texttt{==} смотрит на равенство значений, а \texttt{is} проверяет, один это объект или нет (адреса памяти)». [cite: 105]
			\item[\textcolor{primary}{Дополнение из учебника:}] Различие между сравнением по значению и по идентификатору — ключевой инженерный нюанс Python. 
			\item[\textcolor{primary}{Резюме:}] \texttt{==} проверяет данные, а \texttt{is} проверяет идентичность ссылок. [cite: 107]
		\end{description}
		
		\textcolor{extra_info}{\textbf{Пример на Python:}}
		\begin{lstlisting}
			x = [1, 2, 3]
			y = [1, 2, 3]
			print(x == y) # True (значения одинаковы)
			print(x is y) # False (разные области памяти)
		\end{lstlisting}
		
		\item \textbf{Как устроен \texttt{dict} в Python и что такое коллизии.}
		\begin{description}[leftmargin=1em]
			\item[\textcolor{primary}{Цитата из видео:}] «Dict — это хэш-мапа. Асимптотика в среднем $O(1)$. Коллизия — когда для разных ключей получается один индекс. В Python используется открытая адресация». [cite: 108, 109, 110]
			\item[\textcolor{primary}{Резюме:}] Словарь — хэш-таблица на открытой адресации. [cite: 111, 112] Массив расширяется в 2 раза при заполнении на 66\%. [cite: 113]
		\end{description}
		
		\item \textbf{Что такое итераторы и зачем они нужны.}
		\begin{description}[leftmargin=1em]
			\item[\textcolor{primary}{Цитата из видео:}] «Итератор — объект, позволяющий обойти коллекцию, не держа её в памяти. Обязательны \texttt{\_\_iter\_\_} и \texttt{\_\_next\_\_}». [cite: 114, 115]
			\item[\textcolor{primary}{Дополнение из учебника:}] Итераторы реализуют концепцию "ленивых" вычислений. [cite: 116] Это критично для Big Data. [cite: 117]
			\item[\textcolor{primary}{Резюме:}] Итераторы экономят оперативную память, вычисляя элементы по одному "на лету". [cite: 118]
		\end{description}
		
		\item \textbf{Теория вероятностей: задача на сумму 11.}
		\begin{description}[leftmargin=1em]
			\item[\textcolor{primary}{Цитата из видео:}] «Какая вероятность того, что при подбрасывании двух шестигранных кубиков выпадет сумма 11?». [cite: 119]
			\item[\textcolor{primary}{Резюме:}] Общее количество исходов: $6 \times 6 = 36$. [cite: 120] Благоприятных исхода только два: (5, 6) и (6, 5). [cite: 121] Вероятность $P = 2/36 = 1/18$.
		\end{description}
		
		\item \textbf{Теорема Байеса: формула и применение.}
		\begin{description}[leftmargin=1em]
			\item[\textcolor{primary}{Цитата из видео:}] «Почти всегда выпадает задача на теорему Байеса. Просто заучите формулу». [cite: 122, 123]
			\item[\textcolor{primary}{Дополнение из учебника:}] Формула:
			\[ P(A|B) = \frac{P(B|A)P(A)}{P(B)} \]
			Формула корректирует априорные ожидания с учетом новых фактов. [cite: 124]
			\item[\textcolor{primary}{Резюме:}] Переоценивает вероятности гипотез при поступлении новых данных. [cite: 125]
		\end{description}
		
		\item \textbf{Что такое p-value.}
		\begin{description}[leftmargin=1em]
			\item[\textcolor{primary}{Цитата из видео:}] «p-value — это вероятность получить наблюдаемое различие или более экстремальное, если нулевая гипотеза верна». [cite: 126]
			\item[\textcolor{primary}{Дополнение из учебника:}] Оценивает противоречие данных нулевой гипотезе в A/B тестах. [cite: 127]
			\item[\textcolor{primary}{Резюме:}] Это метрика "удивления". Чем меньше p-value, тем маловероятнее получить результаты случайно. [cite: 128]
		\end{description}
		
		\item \textbf{Статтесты.}
		\begin{description}[leftmargin=1em]
			\item[\textcolor{primary}{Цитата из видео:}] «t-test для нормального распределения, Манн-Уитни — если есть выбросы, Bootstrap — универсально». [cite: 129, 130]
			\item[\textcolor{primary}{Резюме:}] Выбор теста зависит от распределения: параметрические для средних, непараметрические для медиан. [cite: 131, 132]
		\end{description}
		
		\item \textbf{Функции активации и Dropout.}
		\begin{description}[leftmargin=1em]
			\item[\textcolor{primary}{Цитата из видео:}] «Активации добавляют нелинейность. Dropout борется с переобучением, отключая нейроны на трейне». [cite: 133, 134]
			\item[\textcolor{primary}{Дополнение из учебника:}] На инференсе выходы масштабируются на $(1-p)$. [cite: 135, 136]
			\item[\textcolor{primary}{Резюме:}] Активации ломают линейность слоев, а Dropout обучает ансамбль подсетей для повышения устойчивости.
		\end{description}
	\end{enumerate}
	
\end{document}